Table~\ref{tab:nts-factorial} compares DA with Nominate-then-Select (NtS) under no retrieval, global Top-5 retrieval, and parent-balanced retrieval. Within each condition, the two policies reuse the same Diagnostician ranking, criterion states, and criterion-compatible set. The reported difference is NtS minus DA.

\begin{table}[htbp]
\centering
\caption{Paired comparison of Direct-Answer (DA) and Nominate-then-Select (NtS) under three retrieval conditions on the internal held-out split ($N=1000$ per condition). Confidence intervals are paired-bootstrap intervals; negative differences favor DA.}
\label{tab:nts-factorial}
\small
\begin{tabular}{|l|c|c|c|}
\hline
Retrieval condition & DA Top-1 & NtS Top-1 & Difference (95\% CI), pp\\
\hline
No retrieval & 51.0\% & 48.2\% & $-2.8$ [$-4.1$, $-1.5$]\\
\hline
Global Top-5 & 52.4\% & 49.5\% & $-2.9$ [$-4.4$, $-1.4$]\\
\hline
Parent-balanced & 51.8\% & 49.0\% & $-2.8$ [$-4.3$, $-1.3$]\\
\hline
\end{tabular}
\end{table}

NtS reduces committed Top-1 by 2.8--2.9 percentage points in all three retrieval conditions, and each paired-bootstrap confidence interval excludes zero. The detailed discordant counts in Appendix~\ref{app:supporting} show more cases correct only under DA than cases correct only under NtS in every condition. The reduction is therefore not explained by one retrieval setting. The interaction contrasts are small and statistically inconclusive.

Because DA and NtS share the same upstream artifacts, this is the clearest test of finalization policy alone. The result shows that replacing the Diagnostician's proposed primary with the highest-ranked diagnosis selected by the evaluated NtS rule does not improve benchmark agreement. Under the evaluated rule, high gold-label inclusion in a compatible set with a median size of six is not sufficiently selective for reliable primary re-selection.

\paragraph{Additional primary-selection methods.}
Table~\ref{tab:internal-selection-results} summarizes the remaining internal methods relative to the validation-selected parent-balanced DA configuration.

\begin{table}[htbp]
\centering
\caption{Additional internal primary-selection methods relative to DA. Confidence intervals are paired-bootstrap intervals for the Top-1 difference; a dash denotes unavailable paired records.}
\label{tab:internal-selection-results}
\small
\renewcommand{\arraystretch}{1.08}
\begin{tabularx}{\textwidth}{|>{\raggedright\arraybackslash}p{3.2cm}|c|c|c|X|}
\hline
Method & Top-1 & $\Delta$ vs DA & 95\% CI & Interpretation\\
\hline
DA & 0.518 & --- & --- & Reference output\\
\hline
Deterministic fusion & 0.514 & $-0.004$ & --- & Recorded-output summary\\
\hline
Pairwise independent & 0.524 & $+0.006$ & [$-0.010$, 0.022] & Complete configuration\\
\hline
Two-view debate & 0.485 & $-0.033$ & [$-0.049$, $-0.018$] & Complete configuration\\
\hline
\end{tabularx}
\renewcommand{\arraystretch}{1.0}
\end{table}

Deterministic fusion is 0.4 percentage points below DA, but a complete paired case-level record was not retained, so no interval or case-level transition analysis is reported. Independent pairwise comparison is 0.6 percentage points above DA, but its confidence interval includes zero. It also comes from a separate full-pipeline execution, so the difference cannot be attributed only to the pairwise selector. Two-view debate is 3.3 percentage points below DA, and its confidence interval excludes zero. Internal self-consistency and confidence-informed self-consistency remain descriptive because only summary-level records were retained.

\paragraph{Answer to RQ4.}
None of the evaluated methods provides a clear and attributable improvement over DA on the internal held-out set. The matched NtS comparison consistently lowers Top-1, deterministic fusion is slightly lower without retained paired records, the small independent-pairwise increase is statistically inconclusive and confounded by a complete rerun, and debate performs below DA. The observed candidate and compatibility headroom therefore was not converted into a reliable gain by the tested selectors.

\section{Model-Scale and Diagnostic Disagreement Analysis}
\label{sec:model-scale-disagreement}

This section reports the internal model-scale and diagnostic-confusion results that contribute to RQ5. External synthetic transfer, diagnostic-scope expansion, and lexical transfer are reported in Chapter~\ref{ch:external}.

\paragraph{Model scale.}
\ThesisFigure{fig_a_size_robustness.pdf}{Held-out Top-1 across Qwen3 model sizes for the Single baseline and HiED--DA. Shaded bands indicate the observed minimum-to-maximum range across the evaluated model sizes and are not confidence intervals.}{fig:size-robustness}

Single Top-1 increases monotonically from 0.143 at 4B to 0.465 at 32B. HiED--DA Top-1 varies from 0.497 to 0.538 without a monotonic size trend, and its genuine Top-3 Accuracy remains between 0.803 and 0.818. These ranges are narrower than the change observed for Single, but they do not establish model-size invariance. They apply only to the evaluated Qwen3 family and fixed study settings.

The smallest-model runs also contained additional structured-output failures. These execution issues are not diagnostic outcomes. Exact counts and run-integrity notes are reported in Appendix~\ref{app:supporting}, Table~\ref{tab:qwen-size-results}.

\paragraph{Diagnostic confusion.}
The first-listed-label confusion analysis has agreement 0.509, compared with the any-gold committed Top-1 Accuracy of 0.518. The difference occurs because the confusion matrix assigns one first-listed projected benchmark parent to each case, whereas the main Top-1 measure counts a match with any projected benchmark gold parent.

\ThesisFigure{fig_d_confusion_heatmap.pdf}{Confusion matrix on the internal held-out set using the first-listed gold parent for each case.}{fig:confusion-thesis-rewrite}

The largest direction is F41$\rightarrow$F32 with 178 cases, compared with 24 cases in the reverse F32$\rightarrow$F41 direction. Other frequent directions are F39$\rightarrow$F32 (43), \emph{Others}$\rightarrow$F32 (37), and \emph{Others}$\rightarrow$F41 (26). These counts show that the disagreements are not distributed symmetrically and that F32 is a frequent destination. Because the matrix uses one first-listed benchmark label, it is descriptive rather than the primary multilabel evaluation. Chapter~\ref{ch:error} examines the candidate ranks, selection pairs, and criterion-state patterns behind these directions.

\paragraph{Internal contribution to RQ5.}
Within the evaluated Qwen3 family, HiED's ranked Top-3 coverage changes little across model sizes, while committed Top-1 varies over a narrower range than the Single baseline. The internal disagreements are also concentrated in a small number of diagnostic directions, especially F41$\rightarrow$F32. These findings define the internal boundary conditions; Chapter~\ref{ch:external} tests whether the broader stage-wise pattern persists under a second synthetic corpus and other diagnostic settings.

\section{Chapter Summary}

The internal results give five main answers. Retrieval improves the Single baseline more clearly than HiED, and parent-balanced retrieval remains the primary HiED configuration because of the validation contract rather than the held-out ranking. TF--IDF with logistic regression is the strongest same-corpus label predictor, while Single and HiED have almost identical committed Top-1 performance. HiED nevertheless exposes a large difference between candidate availability, criterion compatibility, and committed-primary agreement: 272 of 915 checker-eligible cases contain a benchmark gold parent in both the Top-3 and compatible set but commit another diagnosis. The tested primary-selection interventions do not provide a clear attributable improvement over DA. Finally, HiED's Top-3 coverage is comparatively stable across the evaluated Qwen3 sizes, while the remaining disagreements are concentrated toward a small number of diagnostic directions. Chapter~\ref{ch:external} examines whether these patterns persist under external synthetic distribution shift.
